% ============================================================================
% Fichier  : annexes/Ann_B_outils.tex
% Version  : 1.0 -- novembre 2026
% ============================================================================

\chapter{Outils}
\label{ann:outils}

Les outils servent l'audit~; ils ne le remplacent pas. Chacun produit des
pièces qu'il faut ensuite qualifier, coter et documenter. Les outils de
découverte et d'analyse s'utilisent avec l'accord écrit du commanditaire,
sur les seuls systèmes de l'organisation.

\begin{tableaucours}{L{3.6cm}L{4.4cm}L{4.0cm}}{%
    \tableauHeader{Besoin} &
    \tableauHeader{Outil ou méthode} &
    \tableauHeader{Chapitre}
}
Tenir le registre & Tableur structuré selon le Tableau~4.1, ou base
légère~; un modèle figure à l'annexe~\ref{ann:modeles} &
\ref{chap:cartographie} \\
Découvrir les données & Recherche de motifs par expressions régulières
(\texttt{grep}, \texttt{ripgrep}) sur les partages~; lecture des schémas
de bases & \ref{chap:cartographie} \\
Conduire une évaluation d'impact & Logiciel PIA de la CNIL (libre) ou
modèle tableur & \ref{chap:risque} \\
Analyser les flux d'une application & Mandataire d'interception sur un
appareil de test (\texttt{mitmproxy}), journaux DNS, \texttt{Wireshark} &
\ref{chap:transferts} \\
Repérer les traqueurs d'une application mobile & Analyse statique
(\texttt{MobSF}), bases publiques de traqueurs (Exodus Privacy) &
\ref{chap:cartographie} \\
Exploiter des journaux d'accès & Tableur, \texttt{Python} et
\texttt{pandas} pour le rapprochement avec les plannings &
\ref{chap:securite} \\
Mesurer $\mu$, $U$, $H$ & Tableur ou script~; formules du
Chapitre~\ref{chap:audit} & \ref{chap:audit} \\
Comparer des taux de décision & Tableaux croisés par groupe déclaré &
\ref{chap:profilage} \\
Protéger les données de l'audit & Conteneur chiffré, masquage des
identités dans les papiers & \ref{chap:normes} \\
\end{tableaucours}
\vspace{-0.8em}
\begin{center}{\small\itshape Tableau~B.1. Outils de l'audit de protection des données}\end{center}

\begin{boxinfo}
\textbf{Un outil a ses angles morts.} Une recherche de motifs ne voit pas
les images scannées~; un mandataire d'interception ne voit pas les flux
d'un serveur~; une analyse statique ne dit pas ce qu'une bibliothèque
envoie vraiment. L'auditeur écrit, pour chaque outil, ce qu'il a couvert
et ce qu'il n'a pas pu couvrir.
\end{boxinfo}
